\documentclass[11pt]{scrartcl}
\usepackage{graphicx}
\usepackage{indentfirst}
\usepackage[most]{tcolorbox}
\usepackage{xcolor}
\usepackage{asymptote}
\usepackage{subfigure}
\usepackage{hyperref}
\usepackage{kempu}

\title{Irreducibility}
\author{\LARGE Kempu33334}
\date{\large \today}

\begin{document}
\maketitle

\tableofcontents

\newpage

\section{Preface}

In this handout, we discuss polynomial irreducibility, various methods for proving it, and its applications. Note that we will primarily work over the fields $\mathbb{R}$ and $\mathbb{Q}$, and the ring $\mathbb{Z}$, although the concept of irreducibility applies to any field or ring.

\section{Preliminaries}

We start by defining what it means for a polynomial to be irreducible.

\begin{definition*}[Irreducibility]
A nonconstant polynomial $P(x)$ is \emph{irreducible} over a ring $R$ if whenever $P(x)=f(x)g(x)$ with $f$, $g\in R[x]$, one of $f$, $g$ must be a unit in $R[x]$. Conversely, a polynomial is \emph{reducible} if it can be written as $P(x)=f(x)g(x)$ where neither $f(x)$ nor $g(x)$ is a unit in $R[x]$.
\end{definition*}

Here are some examples of irreducibility/reducibility:

\begin{example*}
\begin{itemize}
\item The polynomial $2x^3-2$ is reducible over $\mathbb{R}$, $\mathbb{Q}$, and $\mathbb{Z}$ as $(2x-2)(x^2+x+1)$.
\item The polynomial $x^2-2$ is reducible over $\mathbb{R}$ as $(x-\sqrt{2})(x+\sqrt{2})$, but irreducible over $\mathbb{Q}$ and $\mathbb{Z}$, since there is no factorization involving only rational numbers or integers.
\item The polynomial $x^2+1$ is irreducible over $\mathbb{R}$, $\mathbb{Q}$, and $\mathbb{Z}$, but reducible over $\mathbb{C}$ as $(x-i)(x+i)$.
\end{itemize}
\end{example*}

Notably, by the Fundamental Theorem of Algebra, all polynomials of degree at least $2$ with complex coefficients are reducible over $\mathbb{C}$.

One of the most useful tools for determining reducibility is the Rational Root Theorem (RRT).

\begin{theorem*}[Rational Root Theorem]
Let $P(x) = a_n x^n + a_{n-1} x^{n-1} + \dots + a_1 x + a_0$ be a polynomial with integer coefficients such that $a_n \neq 0$ and $a_0 \neq 0$. If a rational number $\frac{p}{q}$ written in lowest terms is a root of $P(x)$, then:
\begin{itemize}
\item $p$ must be an integer factor of the constant term $a_0$.
\item $q$ must be an integer factor of the leading coefficient $a_n$.
\end{itemize}
\end{theorem*}

A useful lemma follows from this:

\begin{lemma*}
A polynomial $P(x)$ of degree $2$ or $3$ over a field (such as $\mathbb{Q}$ or $\mathbb{R}$) is irreducible if and only if it has no roots in that field. Thus, if the Rational Root Theorem shows a cubic or quadratic has no rational roots, it is automatically irreducible over $\mathbb{Q}$.
\end{lemma*}
\begin{proof}
We prove the result for $\mathbb{Q}$; the proof for $\mathbb{R}$ is nearly identical.

We start with the if direction. Suppose that $P(x)$ has degree at most $3$ and has no roots in the field $\mathbb{Q}$. Assume for contradiction that $P(x)$ is reducible, so that $P(x)=f(x)g(x)$, where $f$, $g\in\mathbb{Q}[x]$ are both non-constant. Since \[\deg P=\deg f+\deg g\le 3,\] at least one of $f$ or $g$ must have degree $1$; otherwise both would have degree at least $2$, giving \[\deg P=\deg f+\deg g\ge 4\] which is a contradiction.

WLOG, let $f$ be linear. Since $f$ has coefficients in $\mathbb{Q}$, its unique root $r$ also belongs to $\mathbb{Q}$. Then \[P(r)=f(r)g(r)=0\]so $r$ is a rational root of $P$, contradicting our assumption that $P$ has no roots in $\mathbb{Q}$. Therefore, $P(x)$ is irreducible.

We finish with the only if direction. Suppose that $P(x)$ is of degree $\leqslant 3$ and is irreducible. Then, if it has a root $r$ in the field $\mathbb{Q}$, then $x-r$ must be a factor of $P(x)$, which contradicts irreducibility, so we are done.
\end{proof}

Here is an example of RRT and the above lemma used for irreducibility:

\begin{example*}
Prove that the polynomial $P(x) = 2x^3 + x - 1$ is irreducible over $\mathbb{Q}$.
\end{example*}

Since $P(x)$ has degree $3$, it is irreducible over $\mathbb{Q}$ if and only if it has no rational roots. According to the Rational Root Theorem, any rational root $\frac{p}{q}$ in lowest terms must satisfy:
\begin{itemize}
\item $p$ is a factor of the constant term $a_0 = -1$, so $p \in \{\pm 1\}$.
\item $q$ is a factor of the leading coefficient $a_3 = 2$, so $q \in \{\pm 1, \pm 2\}$.
\end{itemize}
This gives a finite list of possible rational roots: $\frac{p}{q} \in \left\{\pm 1, \pm \frac{1}{2}\right\}$. We test each candidate by evaluating the polynomial:
\begin{align*}
P(1) &= 2(1)^3 + (1) - 1 = 2 \\
P(-1) &= 2(-1)^3 + (-1) - 1 = -4 \\
P\left(\frac{1}{2}\right) &= 2\left(\frac{1}{2}\right)^3 + \left(\frac{1}{2}\right) - 1 = \frac{1}{4} + \frac{1}{2} - 1 = -\frac{1}{4} \\
P\left(-\frac{1}{2}\right) &= 2\left(-\frac{1}{2}\right)^3 + \left(-\frac{1}{2}\right) - 1 = -\frac{1}{4} - \frac{1}{2} - 1 = -\frac{7}{4}.
\end{align*}
Since none of the possible candidates are roots, $P(x)$ has no rational roots. Because $\deg(P) = 3$, we conclude that $P(x)$ is irreducible over $\mathbb{Q}$.

\begin{remark*}
The RRT test for irreducibility only applies to polynomials of degree $2$ and $3$. A higher-degree polynomial can lack roots and still be reducible. For example, $x^4 + 2x^2 + 1 = (x^2+1)^2$ has no real or rational roots, but it is clearly reducible.
\end{remark*}

\section{Reduction Maps}

A powerful strategy for analyzing polynomials over $\mathbb{Z}$ is to map them into a simpler algebraic setting where arithmetic is finite. We achieve this by projecting the infinite ring of integer coefficients into a finite field using a modular reduction map.

\begin{definition*}[Reduction Map]
Let $p$ be a prime number. The reduction map modulo $p$ is the function $\phi: \mathbb{Z}[x] \mapsto \mathbb{Z}_p[x]$ that replaces each coefficient of a polynomial with its remainder modulo $p$. Explicitly, given a polynomial $P(x) = a_n x^n + a_{n-1} x^{n-1} + \dots + a_0 \in \mathbb{Z}[x]$, its image under $\phi$ is:
\[ \phi(P) = (a_n \bmod p)x^n + (a_{n-1} \bmod p)x^{n-1} + \dots + (a_0 \bmod p) \]
\end{definition*}

The utility of the reduction map relies entirely on its preservation of algebraic structure. Crucially, $\phi$ is a \emph{ring homomorphism}. This means that reducing the components of an algebraic operation yields the same result as performing the operation first and reducing the final output:
\begin{itemize}
\item \textbf{Additive Preservation:} $\phi(f + g) = \phi(f) + \phi(g)$
\item \textbf{Multiplicative Preservation:} $\phi(fg) = \phi(f)\phi(g)$
\end{itemize}

Because $p$ is a prime number, the finite ring $\mathbb{Z}_p$ is a field. A fundamental structural property of any polynomial ring over a field (such as $\mathbb{Z}_p[x]$) is that it forms an \emph{integral domain}. Consequently, $\mathbb{Z}_p[x]$ contains no zero divisors: if the product of two polynomials is zero, at least one of the underlying polynomials must be zero. That is, if $\phi(f)\phi(g) = 0$, then either $\phi(f) = 0$ or $\phi(g) = 0$. 

\subsection{Gauss' Lemma}
Before testing higher-degree polynomials for irreducibility, we need a way to bridge the gap between factoring over the integers $\mathbb{Z}$ and factoring over the rational numbers $\mathbb{Q}$. Gauss's Lemma provides this bridge.

Firstly, we define a primitive polynomial.

\begin{definition*}[Primitive Polynomial]
A polynomial $P(x) = a_n x^n + \dots + a_0 \in \mathbb{Z}[x]$ is called \emph{primitive} if the GCD of all its coefficients is $1$. That is, $\gcd(a_n, a_{n-1}, \dots, a_0) = 1$.
\end{definition*}

We can now state the lemma.

\begin{lemma*}[Gauss's Lemma]
The product of two primitive polynomials is also a primitive polynomial.
\end{lemma*}
\begin{proof}
Let $f(x)$ and $g(x)$ be primitive polynomials in $\mathbb{Z}[x]$. Let $h(x) = f(x)g(x)$ be their product. We want to show that $h(x)$ is also primitive.

We proceed by contradiction. Suppose $h(x)$ is not primitive. By definition, the greatest common divisor of the coefficients of $h(x)$ must be greater than $1$, meaning there exists some prime number $p$ that divides every coefficient of $h(x)$.

Now, consider the reduction map modulo $p$, $\phi: \mathbb{Z}[x] \mapsto \mathbb{Z}_p[x]$. Because $p$ divides every coefficient of $h(x)$, its image under the reduction map is the zero polynomial: $\phi(h) = 0$

Since $\phi$ is a ring homomorphism, it preserves polynomial multiplication. We can rewrite the image of the product as the product of the individual images: \[\phi(f)\phi(g) = \phi(fg) = \phi(h) = 0.\]

Recall that because $p$ is a prime number, the finite field $\mathbb{Z}_p$ ensures that the polynomial ring $\mathbb{Z}_p[x]$ is an integral domain. Therefore, we must have either $\phi(f) = 0$ or $\phi(g) = 0$. WLOG let $\phi(f) = 0$. Then, every coefficient of $f(x)$ reduces to $0 \pmod p$. This means $p$ divides every coefficient of $f(x)$, which directly contradicts our assumption that $f(x)$ is primitive.

Thus, the initial assumption is false, and there exists no such prime $p$, meaning the greatest common divisor of the coefficients of $h(x)$ is $1$, as desired.
\end{proof}

For proving irreducibility, the most practical takeaway is the immediate corollary of this lemma. It guarantees that if a polynomial cannot be broken down using integers, it cannot be broken down using fractions either.

\begin{theorem*}[Gauss's Lemma for Irreducibility]
Let $P(x) \in \mathbb{Z}[x]$ be a primitive integer polynomial. Then, $P(x)$ is irreducible in $\mathbb{Z}[x]$ if and only if it is irreducible in $\mathbb{Q}[x]$.
\end{theorem*}

For example, consider the polynomial $P(x) = x^4 - x + 1$. Suppose we want to show it is irreducible over $\mathbb{Q}[x]$. By Gauss's Lemma, we do not need to worry about complex fractional factors like $x^2 + \frac{1}{2}x + \frac{3}{4}$; we only need to prove that it cannot be split into integer-coefficient polynomials in $\mathbb{Z}[x]$.

Thus, from now on, we work only in $\mathbb{Z}[x]$.

\subsection{Reduction Maps Revisited}

The following theorem allows us to use reduction maps to prove irreducibility.

\begin{theorem*}
Define the reduction map $\phi: \mathbb{Z}[x] \mapsto \mathbb{Z}_p[x]$ for prime $p$ not dividing the leading coefficient of $P(x)$. Then, if $\phi(P)$ is irreducible and $P$ is primitive, then $P$ is also irreducible. In other words, a primitive polynomial $P(x)$ is irreducible over $\mathbb{Z}[x]$ if it is irreducible over $\mathbb{Z}_p[x]$ for prime $p$ where $p$ does not divide the leading coefficient of $P(x)$.
\end{theorem*}
\begin{proof}
Suppose for contradiction that $P(x)$ is reducible, so that $P(x)=g(x)h(x)$ for nonconstant polynomials $g$, $h\in\mathbb{Z}[x]$. Reducing modulo $p$ gives $\phi(P)=\phi(g)\phi(h)$. Since $\phi(P)$ is irreducible in $\mathbb{F}_p[x]$, one of $\phi(g)$ or $\phi(h)$ must be constant.

Assume without loss of generality that $\phi(g)$ is constant. Since $p$ does not divide the leading coefficient of $P(x)$, it also cannot divide the leading coefficients of $g(x)$ and $h(x)$. Hence the degree of $g(x)$ is preserved under reduction modulo $p$, so \[\deg(\phi(g))=\deg(g)>0\] contradicting the fact that $\phi(g)$ is constant.

Therefore, $P(x)$ must be irreducible.
\end{proof}

We can also state a quick theorem regarding polynomials in $\mathbb{Z}_p[x]$.

\begin{lemma*}[Freshman's Dream]
Let $p$ be a prime. Then, $x^{p^n}+y^{p^n} \equiv (x+y)^{p^n} \pmod{p}$ for integers $n\ge 0$.
\end{lemma*}
\begin{proof}
This simply follows from the binomial theorem. Notice that \[(x+y)^{p^n} = \sum_{k=0}^{p^n} \binom{p^n}{k}x^ky^{p^n-k}.\] Notice that $p \mid \tbinom{p^n}{k}$ for $1 \le k \le p^n-1$ (proof left as an exercise; induction). Thus, the only remaining terms would be at $k = 0$, $p^n$, leaving $x^{p^n}$ and $y^{p^n}$.
\end{proof}

\begin{remark*}
The name comes from the generally false ``identity'' $(x+y)^n = x^n+y^n$ that freshman students tend to make, which just so happens to be true in $\mathbb{Z}_p$ and $n = p$.
\end{remark*}

We can now do our first irreducibility problem.

\begin{example*}[Romanian IMO TST 1998/3/3]
Show that for any $n \in \mathbb{N}$ the polynomial $f(x) = (x^2 +x)^{2^n}+1$ is irreducible over $\mathbb{Z}[x]$.
\end{example*}

Suppose that $f(x)=(x^2+x)^{2^n}+1=g(x)h(x)$ is reducible. Note that $f(x)$ is primitive, so taking the reduction map $\phi: \mathbb{Z}[x] \mapsto \mathbb{Z}_2[x]$, we let $f_1 = \phi(f)$. Then, $f_1$ must be reducible in $\mathbb{Z}_2$. By Freshman's Dream, \[f_1(x) = (x^2 +x)^{2^n}+1 = (x^2+x+1)^{2^n}\] in $\mathbb{Z}_2[x]$, and $x^2+x+1$ is irreducible. Define $g_1(x)$, $h_1(x)$ by $\phi(g)$ and $\phi(h)$ respectively.

Since $\mathbb{Z}_2[x]$ is a Unique Factorization Domain (UFD) and $x^2+x+1$ is irreducible, we can let
\begin{align*}
g_1(x)&=(x^2+x+1)^m \\
h_1(x)&=(x^2+x+1)^{2^n-m}
\end{align*}
for $m\not\in \{0, 2^n\}$, since otherwise, we just get constant polynomials.

Now, let $g(x)=g_1(x)+2g_2(x)$, and $h(x)=h_1(x)+2h_2(x)$. Letting $\omega$ be a root of $x^2+x+1$ in $\mathbb{C}$, if $g$, $h$ are nonconstant, then \begin{align*}
g_1(\omega)&=(\omega^2+\omega+1)^m = 0 \\
h_1(\omega)&=(\omega^2+\omega+1)^{2^n-m} = 0.
\end{align*}
Plugging in $x = \omega$ into $f(x) = g(x)h(x)$, we get that \[(\omega^2+\omega)^{2^n}+1 = 2 = 4g_2(\omega)h_2(\omega).\] Furthermore, $g_2(\omega)h_2(\omega)=a\omega^2+b\omega+c$ for some integers $a$, $b$, $c$ since $\omega^k = \omega^{k-3}$. Because $\omega^2 = -\omega-1$, we can WLOG take $a=0$. But, because \[g_2(\omega)h_2(\omega) = b\omega+c = \tfrac{1}{2}\] we get that $b = 0$ by equating components implying $c = \tfrac{1}{2}$, which isn't an integer, contradiction. Thus, we are done.

\section{Eisenstein's Criterion}

We can now move on to the celebrated Eisenstein's Criterion.

\begin{theorem*}[Eisenstein's Criterion]
Let $P(x) = a_n x^n + a_{n-1} x^{n-1} + \dots + a_1 x + a_0$. If there exists a prime $p$ such that
\begin{itemize}
\item $p$ divides all the coefficients except $a_n$ ($p$ divides $a_{n-1}$, $a_{n-2}$, $\dots$, $a_1$, $a_0$).
\item $p$ does not divide the leading coefficient $a_n$.
\item $p^2$ does not divide the constant term $a_0$.
\end{itemize}
Then $P(x)$ is irreducible over $\mathbb{Q}$ and if $P(x)$ is primitive, then it is also irreducible over $\mathbb{Z}$.
\end{theorem*}
\begin{proof}
Consider the reduction map $\phi:\mathbb{Z}[x]\mapsto\mathbb{Z}_p[x]$. Since $p$ divides all coefficients except the leading coefficient, we have \[\phi(P)=a_nx^n.\] Suppose, for contradiction, that $P(x)$ is reducible over $\mathbb{Q}$. By Gauss' Lemma, we may assume that $P(x)=g(x)h(x)$ for nonconstant primitive polynomials $g$, $h\in\mathbb{Z}[x]$.

Reducing modulo $p$, we get \[a_nx^n=\phi(g)\phi(h).\] Since $\mathbb{Z}_p[x]$ is an UFD and $x$ is irreducible, both $\phi(g)$ and $\phi(h)$ must be powers of $x$. Therefore, $\phi(g)=c x^r$, $\phi(h)=d x^s$ for some nonzero constants $c$, $d\in\mathbb{Z}_p$ with $r+s=n$.

Thus all coefficients of $g(x)$ except its leading coefficient are divisible by $p$, and the same is true for $h(x)$. Write \[g(x)=b_rx^r+pG(x),\qquad h(x)=c_sx^s+pH(x)\] where $b_r$, $c_s$ are not divisible by $p$.

Then, the constant term of $P(x)$ is $a_0=g(0)h(0)$. Because all non-leading coefficients of $g$ and $h$ are divisible by $p$, we have $p\mid g(0)$ and $p\mid h(0)$ unless one of $g(0)$, $h(0)$ is zero.

If one of them is zero, then $a_0=0$, contradicting the condition that $p^2\nmid a_0$ (and in particular $p\nmid a_0$ is not required, but $a_0\neq0$).

Hence both $g(0)$ and $h(0)$ are divisible by $p$, so \[p^2\mid g(0)h(0)=a_0\]contradicting the assumption that $p^2\nmid a_0$. Therefore, $P(x)$ is irreducible over $\mathbb{Q}$. If $P(x)$ is primitive, Gauss' Lemma implies that it is also irreducible over $\mathbb{Z}$.
\end{proof}

Here is an application of Eisenstein's Criterion:

\begin{example*}
Let $\Phi_p(x) = x^{p-1}+\dots+x+1$ be the $p$th cyclotomic polynomial where $p$ is a prime. Then, show that $\Phi_p(x)$ is irreducible.
\end{example*}

Clearly, it is equivalent to prove that $\Phi_p(x+1)$ is irreducible. Notice that \[\Phi_p(x) = x^{p-1}+\dots+x+1 = \dfrac{x^p-1}{x-1}\] so \[\Phi_p(x) = \dfrac{(x+1)^p-1}{x} = \sum_{k=0}^{p-1} \binom{p}{k+1}x^k = x^{p-1}+\binom{p}{p-1}x^{p-2}+\cdots+\binom{p}{1}.\] Finally, the prime $p$ divides all the coefficients except the leading one, and $p^2$ does not divide $\tbinom{p}{1} = p$, so applying Eisenstein's Criterion, we get that $\Phi_p(x)$ is irreducible, as desired.

\section{Examples}

We now present a few examples demonstrating the standard techniques.

\begin{example*}[Yufei Zhao]
Let $f(x) = a_nx^n + a_{n-1}x^{n-1} + \cdots + a_1x + a_0$ be a polynomial with integer coefficients, such that $|a_0|$ is prime and $|a_0| > |a_1| + |a_2| + \cdots + |a_n|$. Show that $f(x)$ is irreducible.
\end{example*}

Let $\alpha$ be a complex zero of $f(x)$. If $\alpha \le 1$, then \[a_n\alpha^n+a_{n-1}\alpha^{n-1}+\cdots+a_1\alpha = -a_0\] so via the Triangle Inequality, we have that
\begin{align*}
|a_0| &= |a_n\alpha^n+a_{n-1}\alpha^{n-1}+\cdots+a_1\alpha| \\
&\le |a_n\alpha^n|+|a_{n-1}\alpha^{n-1}|+\cdots+|a_1\alpha| \\
&\le |a_n|+|a_{n-1}|+\cdots+|a_1|
\end{align*}
which is a contradiction. Thus, $\alpha > 1$. Now, suppose for the sake of contradiction that $f$ is reducible, or $f(x) = g(x)h(x)$ for non-constant integer polynomials $g$ and $h$. Then, at $x = 0$, we get that $f(0) = g(0)h(0)$, so one of $|g(0)|$ or $|h(0)|$ is $1$. WLOG suppose that $|g(0)| = 1$. Then, notice that if $\beta_1$, $\beta_2$, $\dots$, $\beta_k$ are the roots of $g$ and $c$ is the leading coefficient of $g$, then \[|\beta_1\beta_2\cdots \beta_k| = \dfrac{1}{c} \le 1\] which contradicts the fact that all the roots of $f$ must have magnitude strictly greater than $1$. Thus, we are done.

\begin{example*}[Cohn's Irreducibility Criterion]
Let $p = \overline{a_na_{n-1}\cdots a_1a_0}$ be a prime number in base $10$ such that $a_n > 1$. Prove that the polynomial $P(x) = a_nx^n + \cdots + a_1x + a_0$ is irreducible.
\end{example*}

Suppose for the sake of contradiction that $P(x) = Q(x)R(x)$ is reducible where $Q$ and $R$ are non-constant integer polynomials. Taking $x = 10$, we get that \[P(10) = p = Q(10)R(10).\] Hence, one of $|Q(10)|$ or $|R(10)|$ must be equal to $1$, and WLOG let it be $|Q(10)|$.

Now, we make a claim.

\begin{claim*}
For a root $\alpha$ of $P(x)$, we have $|\alpha| \le 5.5$.
\end{claim*}
\begin{proof}
It is obviously true for $|\alpha| < 1$, so assume that $|\alpha| > 1$. Since $P(\alpha) = 0$, we can rearrange as \[a_n\alpha^n = -(a_{n-1}\alpha^{n-1}+\cdots+a_0).\] Since $\alpha \neq 0$ (otherwise, $a_0 = 0$), we can divide out by $\alpha^n$ obtaining: \[a_n = -\left(\dfrac{a_{n-1}}{\alpha}+\cdots+\dfrac{a_0}{\alpha^n}\right).\] Taking magnitudes, we get that
\begin{align*}
a_n &= \left|\dfrac{a_{n-1}}{\alpha}+\cdots+\dfrac{a_0}{\alpha^n}\right| \\
&\le \left|\dfrac{a_{n-1}}{\alpha}\right|+\cdots+\left|\dfrac{a_0}{\alpha^n}\right| \\
&= \dfrac{a_{n-1}}{|\alpha|}+\cdots+\dfrac{a_0}{|\alpha|^n} \\
&< \dfrac{9}{|\alpha|-1}
\end{align*}
via the Triangle Inequality. Rearranging, \[|\alpha| < \dfrac{9}{a_n}+1\] and since $a_n > 1$ is an integer, we get that $|\alpha| < 5.5$, as desired.
\end{proof}

Next, we can express $Q(x)$ in terms of its complex roots $\beta_1$, $\beta_2$, $\dots$, $\beta_m$ (where $m = \deg(Q) \ge 1$): \[Q(x)=c\prod _{j=1}^{m}(x-\beta _{j})\] where $c \in \mathbb{Z}$ is the leading coefficient of $Q(x)$, meaning $|c| \ge 1$. Evaluating the absolute value at $x = 10$: \[|Q(10)|=|c|\prod _{j=1}^m|10-\beta _{j}|.\]Using the triangle inequality and the Claim ($|\beta_j| < 5.5$):\[|10-\beta _{j}|\ge 10-|\beta _j|>10-5.5=4.5\]Since $m \ge 1$, the product satisfies: \[|Q(10)|>1\cdot (4.5)^m>1\]This contradicts $|Q(10)| = 1$, so the assumption that $P(x)$ is reducible must be false, proving that $P(x)$ is irreducible.

\section{More Obscure Theorems}

We now present a few more (albeit less known and more specific) theorems that can help prove irreducibility.

There exists a slightly more general version of Eisenstein that seldom appears.

\begin{theorem*}[Extended Eisenstein's Criterion]
Let $P(x) = a_n x^n + a_{n-1} x^{n-1} + \dots + a_1 x + a_0$. If there exists a prime $p$ such that
\begin{itemize}
\item $p$ divides all the coefficients $a_0$, $a_1$, $\dots$, $a_{k-1}$ for positive integer $k$.
\item $p$ does not divide the leading coefficient $a_n$.
\item $p^2$ does not divide the constant term $a_0$.
\end{itemize}
Then $P(x)$ has an irreducible factor of degree at least $k$.
\end{theorem*}

Here is an example of the Extended Eisenstein's Criterion.

\begin{example*}[IMO 1993/1]
Prove that $P(x) = x^n + 5x^{n-1}+3$ is irreducible if $n \ge 2$ is a positive integer.
\end{example*}

Assume for the sake of contradiction that $P(x)$ is reducible.

By the Extended Eisenstein's Criterion, we can take $p = 3$ and note that $3$ does not divide the leading coefficient, $9$ does not divide the constant term, and $3$ divides all the terms from $x^0$ to $x^{n-2}$, so there exists an irreducible factor of degree at least $n-1$. However, this means that there exists a linear term dividing $x^n+5x^{n-1}+3$. Applying RRT, this must be one of $\{1, -1, 3, -3\}$, but none of these can work, so the irreducible factor must be degree $n$, or the polynomial itself, as desired.

\begin{lemma*}[Bound on Polynomial Roots]
Let $P(x) = x^n + a_{n-1}x^{n-1}+ \cdots + a_0$ be a polynomial such that $|a_i| \le M$ for all $i$. Then all roots of $P(x)$ are less than $M + 1$ in absolute value.
\end{lemma*}

By the triangle inequality, \[|P(x)| > |x|^n - |a_{n-1}x^{n-1}| - \cdots - |a_0| > |x|^n - M(|x|^{n-1} + \cdots + 1) > |x|^n - \dfrac{M|x|^n}{|x|-1}.\] So, if $z$ is a zero of $P$, then $|z|^n - \tfrac{M|z|^n}{|z|-1} < 0$, i.e. $|z| < M + 1$, as desired.

\begin{theorem*}[Perron's Criterion]
Let $P(x) = x^n+a_{n-1}x^{n-1}+\cdots+a_0$ be an integer polynomial such that $|a_{n-1}| > 1+|a_{n-2}|+\cdots+|a_0|$ where $a_0 \neq 0$. Then, $P(x)$ is irreducible over $\mathbb{Q}$.
\end{theorem*}
\begin{proof}
The proof is rather long and involved, so we omit it here.
\end{proof}

\newpage

\section{Problems}

With that, we can move on to the exercises. These are arranged in no particular order.

\begin{problem}[1]{}
The $n$th cyclotomic polynomial $\Phi_n(x)$ is given by \[\Phi_n(x) = \prod_{\stackrel{0\le k\le n-1}{\gcd(k, n) = 1}}(x-e^{2\pi i k/n}).\]
\begin{itemize}
\item[(a)] Prove that $\Phi_n(x)$ is an integer polynomial.
\item[(b)] Prove that $\Phi_n(x)$ is irreducible over the rationals.
\end{itemize}
\end{problem}
\begin{problem}[2]{Selmer's theorem}
Prove that $x^n - x - 1$ is irreducible over the rational numbers $\mathbb{Q}$ for all integers $n \ge 2$.
\end{problem}
\begin{problem}[3]{}
Determine all pairs of polynomials $f$, $g \in \mathbb{Z}[x]$, such that $f(g(x)) = x^{2007}+2x+1$.
\end{problem}
\begin{problem}[4]{}
Let $f(x) = x^4 + 6x^2 + 1$. Show that for any prime $p$, $f(x)$ is reducible over $\mathbb{F}_p$, but $f(x)$ is irreducible over $\mathbb{Z}$.
\end{problem}
\begin{problem}[5]{}
Let $p$ be prime. Show that $f(x) = x^{p-1} + 2x^{p-2}+ 3x^{p-3}+ \cdots + (p - 1)x + p$ is irreducible.
\end{problem}
\begin{problem}[6]{Romania TST 2003}
Let $f(x) \in \mathbb{Z}[x]$ be an irreducible monic polynomial with integer coefficients. Suppose that $|f(0)|$ is not a perfect square. Show that $f(x^2)$ is also irreducible.
\end{problem}
\begin{problem}[7]{MOP 2007}
Let $p(x)$ be a polynomial with integer coefficients. Determine if there always exists a positive integer $k$ such that $p(x) - k$ is irreducible.
\end{problem}
\begin{problem}[8]{}
If $a_1$, $a_2$, $\dots$, $a_n$ are different integers, prove that $P(x) = (x - a_1)(x - a_2)\cdots(x - a_n) - 1$ is irreducible.
\end{problem}
\begin{problem}[9]{}
Let $p$ be a prime number. Prove that all roots of the polynomial $x^p - x + p$ have an absolute value less than $\sqrt[p-1]{p}$.
\end{problem}
\begin{problem}[10]{}
Let $m$, $n$ and $a$ be positive integers and let $p < a - 1$ be a prime. Prove that the polynomial \[f(x) = x^m(x - a)^n + p\] is irreducible.
\end{problem}
\begin{problem}[11]{}
Prove that the polynomial $P(x) = x^n+4$ can be written as a product of some two nonconstant polynomials with integer coefficients if and only if $4 \mid n$.
\end{problem}
\begin{problem}[12]{}
Let $P(x) = a_0 + a_1x + \cdots + a_nx^n$, where $0 < a_0 \le a_1 \le \cdots \le a_n$ are real numbers, then any complex zero of the polynomial satisfies $|z| \le 1$.
\end{problem}
\begin{problem}[13]{}
Prove that $P(x) = x^4+x^3+ax^2 + 9$ is irreducible for all positive integers $a$.
\end{problem}
\end{document}